% =====================================================================================
% Section draft (2026-11-24): RealQM against Hartree--Fock / Heitler--London and against DFT.
% Standalone; intended for realqm_molecular_forces.tex before sec:limits (or the book). Numbers:
% runs_cc/dft_compare.out (cc-pVTZ, PySCF), runs_cc/{lih,beh2,bh,chp}_snode*, runs_cc/be2_snode_*_h005.out.
% NOTE: the RealQM row below uses the HF node radius; replace by the self-contained protocol
% (RealQM core + node from measured IE, runs_cc/sc_*.out) when that queue finishes.
% =====================================================================================
\section{RealQM, Hartree--Fock, Heitler--London and density functional theory}
\label{sec:compare}

\subsection{Two pair constructions, two classical methods}

RealQM represents an electron pair in one of two ways. In the \emph{2hom} construction both electrons
share one domain and one density; in the \emph{2split} construction each electron has its own domain,
and the two meet at a free wall. For H$_2$ the 2hom pair reproduces restricted Hartree--Fock (RHF)
to $5\times10^{-5}$\,E$_h$, and the 2split pair reproduces the Heitler--London (HL) valence-bond
description~\cite{hl,hartree,fock}. RealQM thus contains, as its two pair constructions, the two methods
with which quantum chemistry began in 1927--1930.

These methods are therefore the natural references for diagnosis: when RealQM differs from RHF on a
2hom construction, the difference is in what RealQM adds or leaves out, not in the pair itself. This is
how the missing repulsion between electron groups was found. Two closed-shell He atoms bind falsely in
basic RealQM ($-0.55$\,E$_h$ at $R=2$) where RHF repels ($+0.12$), and Be$_2$ with free walls binds
by $0.18$\,E$_h$ where RHF gives no bond. The difference is Pauli exclusion: HF enforces it by
antisymmetry of one wave function, while basic RealQM's free walls charge nothing where same-spin
electrons meet. Restoring it as a wall condition---an evanescent Robin cost
$\partial_n u+\beta u=0$, $\beta=\sqrt{-2\varepsilon}$ from the domain's own eigenvalue, between closed
shells on different atoms, and a node condition on the spherical part of the valence at the core---removes
the false bonds (He$_2$: $+0.056$, $+0.011$, $+0.0025$ at $R=2,3,4$, against RHF $+0.121$, $+0.014$, $+0.001$; Be$_2$ repulsive at every $R$) and
leaves no fitted constant.

Two points follow. First, comparing with methods a hundred years old is not a measure of progress
beyond them; it measures whether RealQM's own construction contains their physics. Second, RealQM is
not HF in disguise: HF under-binds every molecule (missing correlation), while RealQM's errors differ in
sign from molecule to molecule. Its missing pieces are different, and they are located in identifiable
parts of the model---walls and cores---rather than in a global term.

\subsection{RealQM and density functional theory}

Standard quantum chemistry advanced beyond HF in two directions: correlated wave-function methods
(configuration interaction, coupled cluster, of steeply rising cost) and density functional theory
(DFT)~\cite{hk,ks}, which replaced the $3N$-dimensional wave function by the three-dimensional density
and is today by far the most used method. Since RealQM also works in three dimensions, DFT is its
real competitor. Table~\ref{tab:dft} sets RealQM beside HF, three density functionals and a near-exact
reference.

\begin{table}[h]
\centering\small
\caption{Bond length $R_e$ (bohr) and binding energy $D_e$ (E$_h$; BeH$_2$: atomisation; CH$^+\to$C$^+$+H).
HF, LDA, PBE, B3LYP and CCSD(T) in cc-pVTZ (PySCF~\cite{pyscf}), five-point scans. RealQM: Pauli walls with
RealQM's own 2hom 1s$^2$ core and the node radius set once per element by the atom's measured ionisation
energies; axisymmetric, fixed mesh $h=0.05$. Last two rows: mean absolute error against CCSD(T).}
\label{tab:dft}
\begin{tabular}{lcccccc}
\hline
 & CCSD(T) & RealQM & HF & LDA & PBE & B3LYP\\
\hline
LiH $R_e$ & 3.025 & 2.973 & 3.046 & 3.041 & 3.041 & 3.014\\
BeH$_2$ $R_e$ & 2.511 & 2.433 & 2.525 & 2.543 & 2.536 & 2.512\\
BH $R_e$ & 2.325 & 2.334 & 2.316 & 2.381 & 2.372 & 2.338\\
CH$^+$ $R_e$ & 2.129 & 2.251 & 2.109 & 2.199 & 2.184 & 2.155\\
LiH $D_e$ & 0.0908 & 0.059 & 0.054 & 0.097 & 0.085 & 0.093\\
BeH$_2$ $D_e$ & 0.2334 & 0.286 & 0.199 & 0.256 & 0.236 & 0.247\\
BH $D_e$ & 0.1348 & 0.136 & 0.098 & 0.143 & 0.129 & 0.133\\
CH$^+$ $D_e$ & 0.1559 & 0.237 & 0.113 & 0.171 & 0.154 & 0.155\\
\hline
$|\Delta R_e|$ & -- & 2.7\% & 0.6\% & 1.9\% & 1.5\% & 0.5\%\\
$|\Delta D_e|$ & -- & 28\% & 27\% & 8\% & 3\% & 2\%\\
\hline
\end{tabular}
\end{table}

On geometry RealQM is close to the local density approximation (2.7\% against 1.9\%), with BH within
0.4\% and CH$^+$ 6\% long; on energies it is at the level of HF in size (28\% against 27\%) but not in
sign---LiH under-bound, BeH$_2$ and CH$^+$ over-bound, BH right---and well behind the functionals. This is the honest present standing.
Numbers alone, however, do not settle what kind of thing is being compared.

\paragraph{What DFT is a model of.}
The Hohenberg--Kohn theorem~\cite{hk} states that the ground-state energy is a functional of the
density; it proves existence, not form. Kohn and Sham~\cite{ks} compute the kinetic energy of fictitious
non-interacting electrons with the same density and collect everything else into the
exchange--correlation functional $E_{xc}[\rho]$. The exact $E_{xc}$ is defined through the many-electron
wave function~\cite{levy}: evaluating it requires solving the $3N$-dimensional Schr\"odinger equation,
which is the problem DFT was meant to avoid. Every functional in use is therefore a chosen formula, its
constants set by a model system (the uniform electron gas, LDA), by exact constraints (PBE~\cite{pbe}),
or by fitting to molecular data (B3LYP~\cite{b3lyp} and successors with dozens of parameters). An exact
theory that cannot be computed makes no prediction and cannot be tested; what is tested is always one
of these formulas.

The criticism is made inside the field. Medvedev, Perdew and co-workers~\cite{medvedev} found that
functionals developed since about 2000 improved energies while their densities became worse---good
answers, in Hammes-Schiffer's phrase, possibly for the wrong reasons~\cite{hammes}. Cohen, Mori-S\'anchez
and Yang~\cite{cohen} trace the systematic failures (stretched bonds, band gaps, charge transfer) to
delocalisation and static-correlation errors of the approximations, and Burke~\cite{burke} notes that
no fix has become a universal tool of known performance. The Hartree term moreover lets each electron
repel its own density; the approximate $E_{xc}$ cancels this only in part.

What DFT computes in three dimensions is therefore not offered as three-dimensional physics. Its one
physical 3D object is the total density. The Kohn--Sham electrons are declared fictitious, and $E_{xc}$
encodes correlations in $3N$-dimensional configuration space. DFT is a three-dimensional
\emph{computational method} for a theory posed in $3N$ dimensions.

\paragraph{What RealQM is a model of.}
RealQM is a model of three-dimensional physics by construction: each electron, or pair, is a real
charge density occupying its own region of space; it obeys a Schr\"odinger equation in that region;
electrons interact by Coulomb's law between distinct densities only, so that no electron repels
itself; and Pauli exclusion is spatial---no two same-spin electrons share space. Every object in the
model is offered as physical, and every approximation is a statement about how electrons behave, which
can be tested on its own. The false bonds of He$_2$ and Be$_2$ located one such statement (``a free wall
charges nothing where same-spin electrons meet'') and showed it to be wrong.

The comparison with DFT is thus not between two approximations to the same thing. DFT approximates an
unknown functional of an uncomputable theory; RealQM is a computable model whose rules are physical
claims. Measured by the same standard---a fixed model, its parameters set once from observation, its
predictions compared with experiment on molecules not used to set them---RealQM in its basic form has
no parameters, and its extension with Pauli walls adds one per element (the node radius, set by the
atom's measured ionisation energy). Its present accuracy on energies is below that of non-empirical
DFT. Its case rests on two things DFT does not offer: a physical picture in which a failure points to
the missing physics, and a model whose objects can be computed and interpreted in three dimensions as
they stand.

\paragraph{What remains to be shown.}
RealQM keeps this standing only if each added rule is itself three-dimensional physics. The
evanescent wall cost is: it is the energy of the tail a wall cuts off, set by the domain's own
eigenvalue. The core node is not yet: it is justified by orthogonality to the 1s orbital, a statement
of standard quantum mechanics about overlapping orbitals, which RealQM's disjoint domains do not have.
A three-dimensional account of how far, and at what cost, valence charge penetrates a core is the open
problem that decides whether the extended model is RealQM's own or partly borrowed.

% ---- bibitems to add -----------------------------------------------------------------
% \bibitem{hl} W.\ Heitler and F.\ London, Wechselwirkung neutraler Atome und hom\"oopolare Bindung nach der
%   Quantenmechanik, \emph{Z.\ Phys.}\ \textbf{44}, 455 (1927).
% \bibitem{hartree} D.\ R.\ Hartree, The wave mechanics of an atom with a non-Coulomb central field,
%   \emph{Proc.\ Camb.\ Phil.\ Soc.}\ \textbf{24}, 89 (1928).
% \bibitem{fock} V.\ Fock, N\"aherungsmethode zur L\"osung des quantenmechanischen Mehrk\"orperproblems,
%   \emph{Z.\ Phys.}\ \textbf{61}, 126 (1930).
% \bibitem{hk} P.\ Hohenberg and W.\ Kohn, Inhomogeneous electron gas, \emph{Phys.\ Rev.}\ \textbf{136}, B864 (1964).
% \bibitem{ks} W.\ Kohn and L.\ J.\ Sham, Self-consistent equations including exchange and correlation effects,
%   \emph{Phys.\ Rev.}\ \textbf{140}, A1133 (1965).
% \bibitem{levy} M.\ Levy, Universal variational functionals of electron densities, \emph{Proc.\ Natl.\ Acad.\ Sci.\ USA}
%   \textbf{76}, 6062 (1979).
% \bibitem{pbe} J.\ P.\ Perdew, K.\ Burke and M.\ Ernzerhof, Generalized gradient approximation made simple,
%   \emph{Phys.\ Rev.\ Lett.}\ \textbf{77}, 3865 (1996).
% \bibitem{b3lyp} A.\ D.\ Becke, Density-functional thermochemistry. III. The role of exact exchange,
%   \emph{J.\ Chem.\ Phys.}\ \textbf{98}, 5648 (1993).
% \bibitem{medvedev} M.\ G.\ Medvedev, I.\ S.\ Bushmarinov, J.\ Sun, J.\ P.\ Perdew and K.\ A.\ Lyssenko, Density
%   functional theory is straying from the path toward the exact functional, \emph{Science} \textbf{355}, 49 (2017).
% \bibitem{hammes} S.\ Hammes-Schiffer, A conundrum for density functional theory, \emph{Science} \textbf{355}, 28 (2017).
% \bibitem{cohen} A.\ J.\ Cohen, P.\ Mori-S\'anchez and W.\ Yang, Insights into current limitations of density
%   functional theory, \emph{Science} \textbf{321}, 792 (2008).
% \bibitem{burke} K.\ Burke, Perspective on density functional theory, \emph{J.\ Chem.\ Phys.}\ \textbf{136}, 150901 (2012).
% (pyscf already in the article.)  Check page numbers of hl/hartree/fock/levy/b3lyp against the originals.
